% Part: incompleteness
% Chapter: theories-computability
% Section: extensions-of-q-not-decidable

\documentclass[../../../include/open-logic-section]{subfiles}

\begin{document}

\olfileid{inc}{tcp}{cqn}

\olsection{$\Th{Q}$નાં સુસંગત વિસ્તરણો અનિર્ણેય છે}

\begin{explain}
યાદ રાખો કે સિદ્ધાંત
\emph{સુસંગત} ત્યારે કહેવાય
જ્યારે તે કોઈપણ
સૂત્ર~$!A$ માટે
$!A$ અને $\lnot !A$
બંને સાબિત ન કરે.
વિરોધાભાસમાંથી કંઈપણ તારવી
શકાય છે, એટલે અસુસંગત
સિદ્ધાંત તુચ્છ હોય છે:
દરેક વાક્ય સાબિત થાય છે.
સ્પષ્ટ છે કે સિદ્ધાંત
$\omega$-સુસંગત હોય તો
સુસંગત પણ હોય. પરંતુ
સુસંગતતા નબળી શરત છે
(એટલે સુસંગત છતાં
$\omega$-અસુસંગત સિદ્ધાંતો
હોય છે). વધુ મજબૂત પ્રમેય
મેળવવા
\olref[oqn]{thm:oconsis-q}માં
લેવાયેલી ધારણાને સામાન્ય
સુસંગતતા સુધી નબળી કરી
શકીએ.
\end{explain}

\begin{lem}
``સાર્વત્રિક સંગણનીય સંબંધ''
હોતો નથી. એટલે એવો કોઈ
દ્વિચલ સંગણનીય સંબંધ
$R(x,y)$ નથી કે જે નીચેનો
ગુણધર્મ ધરાવે: $S(y)$
કોઈપણ એકચલ સંગણનીય
સંબંધ હોય, તો એવું~$k$
હોય કે દરેક~$y$ માટે
$S(y)$ સાચું ત્યારે અને
માત્ર ત્યારે જ $R(k,y)$
સાચું હોય.
\end{lem}

\begin{proof}
માનો કે $R(x,y)$ સાર્વત્રિક
સંગણનીય સંબંધ છે.
$S(y)$ તરીકે $\lnot R(y,y)$
સંબંધ લો. $S(y)$ સંગણનીય
હોવાથી કોઈ~$k$ માટે
$S(y)$ એ $R(k,y)$ની
સમકક્ષ છે. પરંતુ ત્યારે
$S(k)$ એ $R(k,k)$ અને
$\lnot R(k,k)$ બંનેની
સમકક્ષ થાય છે; આ
વિરોધાભાસ છે.
\end{proof}


\begin{thm}
$\Th{T}$ એ $\Th{Q}$ને
સમાવતો કોઈપણ સુસંગત
સિદ્ધાંત હોય. તો $\Th{T}$
નિર્ણેય નથી.
\end{thm}


\begin{proof}
માનો કે $\Th{T}$ એ
$\Th{Q}$નું સુસંગત,
નિર્ણેય વિસ્તરણ છે.
$\Th{T}$ વડે સાર્વત્રિક
સંગણનીય સંબંધ વ્યાખ્યાયિત
કરીને વિરોધાભાસ મેળવીશું.

$R(x,y)$ નીચેની શરત હોય
ત્યારે અને માત્ર ત્યારે જ
સાચો લો:
\begin{quote}
$x$ એ સૂત્ર $!D(u)$નો
સંકેત છે અને $\Th{T}$માં
$!D(\num y)$ સાબિત થાય છે.
\end{quote}
$\Th{T}$ નિર્ણેય છે એવી
ધારણા હોવાથી $R$ સંગણનીય
છે. હવે $R$ સાર્વત્રિક છે
તે બતાવીએ. $S(y)$ કોઈપણ
સંગણનીય સંબંધ હોય, તો
તેને $\Th{Q}$માં, એટલે
$\Th{T}$માં પણ,
$!D_S(u)$ સૂત્ર નિરૂપે છે.
તેથી દરેક~$n$ માટે
\begin{eqnarray*}
S(\num n) & \lif & T \vdash !D_S(\num n) \\
& \lif & R(\Gn{!D_S(u)}, n)
\end{eqnarray*}
અને
\begin{eqnarray*}
\lnot S(\num n) & \lif & T \vdash \lnot !D_S(\num n) \\
& \lif & T \not\vdash !D_S(\num n) \quad \text{(કારણ કે $\Th{T}$ સુસંગત છે)} \\
& \lif & \lnot R(\Gn{!D_S(u)}, n).
\end{eqnarray*}
એટલે દરેક~$y$ માટે
$S(y)$ સાચું ત્યારે અને
માત્ર ત્યારે જ
$R(\Gn{!D_S(u)}, y)$
સાચું છે. તેથી $R$
સાર્વત્રિક છે, અને માંગેલો
વિરોધાભાસ મળી ગયો.
\end{proof}

``સાચું અંકગણિત'' એટલે
સિદ્ધાંત~$\Setabs{!A}{ \Sat{\Nat}{!A}}$,
અર્થાત્ પ્રમાણભૂત અર્થઘટનમાં
સાચાં અંકગણિતીય વાક્યોનો ગણ.

\begin{cor}
સાચું અંકગણિત નિર્ણેય નથી.
\end{cor}

%In Section~\olref{undefinability:truth:section} we will state a stronger
%result, due to Tarski.

\end{document}
